\documentclass[a4paper,oneside,12pt]{amsart}

\usepackage[spanish]{babel}
\usepackage[utf8]{inputenc}
%\usepackage[latin1]{inputenc}
\usepackage{latexsym}
\usepackage{multicol}
\usepackage{enumerate}
\usepackage{enumitem}
\usepackage[heightrounded]{geometry}
\geometry{top=2cm,bottom=2cm,left=2cm,right=2cm}
\usepackage{graphicx}
\usepackage{tabularx}
\usepackage[spanish]{cleveref}
\usepackage{printsudoku}
\usepackage{tikz}
\usetikzlibrary{patterns} % Load the patterns library
\usepackage{caption}
\usepackage{subcaption}
\usepackage{booktabs}

\DeclareMathOperator{\cond}{cond}
\DeclareMathOperator{\Id}{Id}
\DeclareMathOperator{\re}{Re}
%\DeclareMathOperator{\im}{Im}
\newcommand{\I}{\mathbb{I}}
\newcommand{\R}{\mathbb{R}}
\newcommand{\Q}{\mathbb{Q}}
\newcommand{\C}{\mathbb{C}}
\newcommand{\Z}{\mathbb{Z}}
\newcommand{\N}{\mathbb{N}}
\DeclareMathOperator{\im}{Im}

\newcommand{\nombremateria}
{Investigaci\'on Operativa (M) \\ Introducción a la Investigación Operativa y Optimización (LCD)}
\newcommand{\nombrecuatrimestre}
    {Segundo cuatrimestre de 2025}
\newcommand{\numeroparcial}
    { Recuperaotrio del Primer Parcial }
\newcommand{\fecha}
    {9 de Diciembre de 2025}
\newcommand{\numerotema}
    {\bf 1}

%\oddsidemargin -1cm \evensidemargin -1cm \textwidth 17.5cm
%topmargin 1.2cm \textheight 25cm
%\setlength{\textheight}{25cm}

\def \ds{\displaystyle}
\theoremstyle{definition}
\newtheorem{quest}{Ejercicio}
\thispagestyle{empty}

\begin{document}

    \begin{minipage}[t]{0.55\textwidth}
    \begin{tabularx}{\textwidth}{Xl}
        Nombre: & LU: \hphantom{5cm}\\[0.5em]
        Nombre: & LU: \\[0.5em]
        Nombre: & LU: \\[0.5em]
    \end{tabularx}
    \end{minipage}
    \hfill
    \begin{minipage}[t]{0.35\textwidth}
        \begin{tabular}[c]{|c|c|c|}
            \hline
            1 & 2 & Calificación \\ \hline \hline
            \hspace{1.2 cm} & \hspace{1.2 cm} & \hspace{1.5 cm} \\
            \hspace{1.2 cm} & \hspace{1.2 cm} & \hspace{1.5 cm} \\ \hline
        \end{tabular}
    \end{minipage}
%\hfill \hfill

%\vspace{-1cm}
%
%\noindent Nombre y apellido:
%
%\vspace{0.2cm}
%
%\noindent Número de libreta:

%\noindent
\hrulefill

\smallskip
\renewcommand{\baselinestretch}{1.1}

\begin{center}
	\large \textbf{\nombremateria} 
	
	\numeroparcial -- \fecha
\end{center}

\vspace{.6cm}
\renewcommand{\theenumi}{\textbf{\arabic{enumi}}}

\begin{center}
    \emph{Entregar el Ejercicio 1 y el Ejercicio 2 en hojas separadas.}
\end{center}
%%% EMPIEZA EL EJERCICIO 1 %%%

\begin{quest} 
    \texttt{[5 pts.]} El juego Sudoku consiste en completar una cuadrícula de $9\times 9$ con números del $1$ al $9$ de
    manera tal que no se repita ningún número en cada columna, en cada fila y en cada uno de los nueve bloques de
    $3\times 3$. En cada instancia, algunas casillas vienen completadas.
    \begin{enumerate}[label=\alph*)]
        \item \verb|[2.5 pts]| Elaborar un modelo de Programación Lineal Entera para resolver el Sudoku, respetando
        las casillas que ya vienen completadas.
        \item \verb|[1.5 pts]| Implementar el modelo en SCIP para resolver la siguiente instancia de Sudoku:
        \begin{figure}[h]
            \centering
            \cluefont{\normalsize}\cellsize{1.5\baselineskip}
            \sudoku{sudoku.txt}
            \caption{Sudoku}\label{fig:ans1}
        \end{figure}

        Interpretando el resultado obtenido, completar la instancia de Sudoku en esta hoja.
        \item \verb|[1 pt]| Utilizando SCIP, decidir si la solución obtenida en el ítem anterior es única.
    \end{enumerate}



\end{quest}

%%% TERMINA EL EJERCICIO 1 %%%

%%% EMPIEZA EL EJERCICIO 2 %%%

\begin{quest}
    \texttt{[5 pts.]} Una empresa de logística quiere construir un modelo lineal para predecir el tiempo de entrega
    $t_i$ (en horas) a partir de varias características de cada envío: $a_i$ la distancia en km., $b_i$ la cantidad de
    paradas intermedias, $c_i$ si el envío es nocturno (0/1) y $d_i$ la cantidad de objetos frágiles en el envío. Se
    desea ajustar el siguiente modelo:
    \[t_i \approx \beta_0 + \beta_1 a_i + \beta_2 b_i + \beta_3 c_i + \beta_4 d_i\]
    con el objetivo de minimizar el error absoluto medio:
    \[\frac{1}{n}\sum_{i=1}^n |t_i - (\beta_0 + \beta_1 a_i + \beta_2 b_i + \beta_3 c_i + \beta_4 d_i)|\]
    sujeto a ciertas restricciones propias del problema:
    \begin{enumerate}
        \item $\beta_1\geq 0$, $\beta_2\geq 0$ y $\beta_4\geq 0$ pues el tiempo de entrega debe aumentar
        proporcionalmente a la distancia recorrida, la cantidad de paradas intermedias y la cantidad de objetos
        frágiles que se entregan. Por otro lado, como
        durante la noche hay menos tránsito, debe valer que $\beta_3 \leq 0$.\footnote{En SCIP, al definir la variable
            $\beta_3$, se pueden ingresar los argumentos \texttt{ub=0, lb=None} para determinar que es no positiva.}
        \item para mejorar la interpretabilidad, a lo sumo tres de los coeficientes $\beta_1, \beta_2, \beta_3, \beta_4$
        pueden ser no nulos.
        \item para garantizar precisión, se busca que al menos el $75\%$ de las observaciones tengan un error
        absoluto menor o igual a una hora.
    \end{enumerate}

    Teniendo en cuenta la formulación del problema:
    \begin{enumerate}[label=\alph*)]
        \item \verb|[3 pts.]| proponer un modelo de Programación Lineal Entera Mixta que, dados un conjunto de datos
        $\{(t_i, a_i,b_i,c_i,d_i)\}_{i=1}^n$, calcule los coeficientes óptimos para la regresión lineal considerando
        las restricciones.
        \item \verb|[2 pts.]| utilizando SCIP, resolver el modelo para hallar los coeficientes óptimos para el
        conjunto de datos
        proporcionados en el notebook. Escribir en la hoja el valor de la función objetivo en el óptimo y los valores
        hallados para $\beta_0, \beta_1,\beta_2,\beta_3,\beta_4$.
    \end{enumerate}
\end{quest}

%%% TERMINA EL EJERCICIO 2 %%%

%%% TERMINA EL PARCIAL %%%



\end{document}

